% Official SCAR LaTeX Template
% Written by Karl Welzel
% Last updated on 21 January 2024
% This document is published under a CC0 1.0 Universal License
% To view a copy of this license, visit http://creativecommons.org/publicdomain/zero/1.0


% STEP 1: Choose your area
% Replace the string after the equality sign with one of "socialsciences", 
% "humanities" or "mpls" to automatically use the correct citation style.
\documentclass[area=socialsciences,review]{scararticle}
% Ignore any warning "Unused global option(s): [area=socialsciences]"


% STEP 2: Add custom LaTeX packages
% If any special LaTeX packages or settings are required, add them here.


% STEP 3: Adapt the following settings
% - title: The full title of your paper
% - shorttitle: The shortened version of your paper title
%   - This needs to fit in the header. You may repeat the full title if it fits.
% - subtitle: The subtitle of your paper (optional, discouraged)
% - author: The names of all authors delimited by commas (used for metadata)
% - shortauthor: The shortened names of the authors
%   - Shorten each name as "Angela Dorothea Merkel" -> "Merkel, A. D."
%   - Delimit with commas and "and"
%   - For four or more authors shorten to "Merkel, A. D., et al."
% - authorline: The names of authors delimited by commas
%   - Add a \textsuperscript{i} after each comma such that each affiliation gets
%     its own number
%   - To acknowledge financial support from other institutions use an
%     unnumbered "Acknowledgements" section at the end of the article
% - affiliation: The affiliations in the order given in "authorline"
%   - Set each affiliation on its own line and start each new line by "\item"
% - editor: leave blank
% - reviewer: leave blank
% - copyeditor: leave blank
% - receiveddate: leave blank
% - publisheddate: leave blank

\scarinfo{
    title={Preparation of Papers for St Catherine's Academic Review},
    shorttitle={Preparation of Papers for SCAR},
    author={First Author, Second Author},
    shortauthor={Author, F. and Author, S.},
    authorline={First Author\textsuperscript{1} and Second Author\textsuperscript{2}},
    affiliation={
        \item National Institute of Standards and Technology (\email{author@boulder.nist.gov})
        \item Department of Physics, Colorado State University (\email{author@lamar.colostate.edu})
    },
    editor={},
    reviewer={},
    copyeditor={},
    receiveddate={},
    publisheddate={},
}

\scarheaders{}

\setcounter{page}{1}

% STEP 4: Supply your own bibliography file
% Export all of your references in BibLaTeX (or BibTeX) format and save them in
% the "bibliography.bib" file.
\addbibresource{bibliography.bib}

\begin{document}

\maketitle

% STEP 5: Include abstract and keywords
% Summarize your work with an abstract. Only use one paragraph.
% Supply between 3 and 7 keywords related to your research.
\begin{scarabstractandkeywords}%
    This document outlines formatting guidance and instructions for papers which will appear in St. Catherine’s Academic Review. The Abstract should summarize the contents of the paper and should consist of one paragraph only. It is followed by three to seven keywords.
    \begin{scarkeywords}%
        One,
        Two,
        Three,
        Four,
        Five,
        Six,
        Seven.
    \end{scarkeywords}
\end{scarabstractandkeywords}


% STEP 6: Write the main content
% Consult the text below for some more guidance on how the paper should be
% written. It also includes some examples on how to use citations, figures and
% tables in LaTeX. Happy writing!

\section{Style}
Papers to be submitted to SCAR must be prepared according to the instructions presented here.
They can be prepared using either Microsoft Word or \LaTeX{}.
This is the \LaTeX{} template.
All papers will be typeset with LaTeX for publication.
By using this template and following the style suggestions, you make the work of the production team significantly easier.

\subsection{Title and Authors}
Follow the \LaTeX{} comments for the \verb|\scarinfo| command at the top of this document.
The \verb|shorttitle| and \verb|shortauthor| fields are used in the headers and need to be short enough to avoid overflow.
The affiliations of all authors must include the relevant institute and the university.
At least one author must list their institutional email address.

\subsection{Structure}

As exemplified above the paper should be structured using sections and subsections (and subsubsections if necessary).
These are always numbered, and the headings should be set in title case.\footnote{Footnotes may be used to supply additional information but should generally be used sparingly.}

\subsubsection{Subsubsection heading}
This is a subsubsection.

\section{General formatting}

\subsection{Spelling}
The use of British English spelling according to the Oxford English Dictionary is encouraged. Alternatively, American spelling is also accepted if it is used consistently.

\subsection{Citations}
This \LaTeX{} document uses BibLaTeX for bibliography management.
As described in the \LaTeX{} comments above you should export your references in BibLaTeX format and save it as \verb|bibliography.bib|.
The citation style is automatically adapted to the chosen research area.
For Social Sciences it is “APA”, for Humanities it is “MLA” and for MPLS articles it is the \verb|numeric| style.
Make sure to \emph{always} include the DOI for each journal article that you cite.
This makes it easier for readers to find the exact version you are referring to.
Citations can both be used within a sentence (like “\textcite[p. 12]{erdHos1960evolution} claim that …”) and at the end of a claim (like “… is the topic of ongoing research \parencite{faust_microbial_2012}”).
Additional information on how to use the BibLaTeX citation commands can be found here: \url{https://www.overleaf.com/learn/latex/Articles/Getting_started_with_BibLaTeX#Citation_Commands}.

\subsection{Figures}
Figures should be provided as high-resolution images and contain as little text as possible.
The number and caption appear below the centred figure.
The caption is lower case except for proper nouns and the first word of a sentence.
Figures are numbered in sequence with sections.
\Cref{fig:scar_logo} is given as an example.
You may use colour, but all figures should be legible in black and white.
Do not rely on the specific placement of a figure. Instead refer to it using its number.
This is because \LaTeX{} attempts to find the placement that is best for the text flow and so might move a figure from the current position to the top of the page or even to the following page.

\begin{figure}[htbp]
    \centering
    \includegraphics[scale=0.2]{scarlogo.png}
    \caption{Note how the label of this caption is 2.1, in sequence with the section.}
    \label{fig:scar_logo}
\end{figure}

\subsection{Tables}
Tables should not use any vertical lines and exactly three horizontal lines.
One at the top (\verb|\toprule|), one after the headings (\verb|\midrule|) and one at the bottom (\verb|\bottomrule|). 
The title and number appear before the table. 
The caption is lower case except for proper nouns and the first word of a sentence. 
Tables are also numbered in sequence with sections.
\Cref{tab:example_label} is given as an example.
More information on how to use tables in \LaTeX{} can be found here: \url{https://www.overleaf.com/learn/latex/Tables}.
If necessary, tables can be moved to their own page and rotated to allow for more or wider columns.
This can be achieved using the \verb|sidewaystable| environment of the \verb|rotating| package.

\begin{table}
    \centering
    \caption{An example table caption}
    \begin{tabular}{llp{25em}}
        \toprule
        Header A & Header B & Column with linewraps \\
        \midrule
        Entry A & Entry B & \lipsum[2] \\
        Entry C & Entry D & \lipsum[3] \\
        \bottomrule
    \end{tabular}
    \label{tab:example_label}
\end{table}

% Add your acknowledgements here
\section*{Acknowledgements}
Use an unnumbered first level heading. You are required to declare any funding and competing interests.

\printbibliography

\end{document}
